% MT004 / explanation — canonical editable content.
先区分两种“参数”：\(w,b\)是模型从数据中学出来的权重和偏置；\(C\)是训练前指定的超参数，控制拟合与正则化的权衡。交叉验证比较的是不同\(C\)的效果，每次比较仍需要重新学习一组\(w,b\)。


补充例子：共有100条训练记录，做五折CV，把它们分成5组，每组20条。对一个候选\(C\)，每次用80条训练、20条验证，轮流做5次，再平均验证成绩。更换\(C\)后重复比较。最终测试集始终单独留着。



\begin{itemize}

\item \textbf{A正确。} 五个验证折来自原训练池。某条记录在这一轮用于验证，下一轮可以参与训练；它不会同时参与同一轮的拟合和验证。

\item \textbf{B错误。} 用验证成绩选\(C\)。若不断根据测试成绩修改\(C\)，测试集已经参与了模型选择，最后的成绩就不再是独立评价。

\item \textbf{C正确。} 选定\(C\)后，用全部100条训练记录重训最终模型。CV中的每个模型只见过80条，最终模型见过100条；即使\(C\)相同，学出的\(w,b\)也可能不同。

\item \textbf{D正确。} 例如只试\(C=0.01\)和\(C=100\)，可能漏掉中间更合适的强度。因此搜索间隔过大“可能”导致表现差，但不能仅凭测试差就确定是这个原因。

\end{itemize}


若还有标准化、词表选择等步骤，也要在每一折的80条训练记录上拟合，再应用到其20条验证记录。否则验证记录的信息会提前进入模型。


\textbf{出处：} 2020B期中Quiz Q4，原卷第2–3页（5分）


\textbf{答案依据：} 2020B期中 Q4 随卷／配套参考答案；条件与解释经复核。
